\section{Conclusion}
\label{sec:conclusion}
In this work, we study benchmark contamination in the context of large language models and evaluate existing decontamination methods.
We show that existing detection methods can not detect test cases with simple variations. 
We demonstrate that if such variation of test data is not eliminated, a 13B model can easily overfit the test benchmark and achieve drastically high performance.
To address this, we propose a new detection method LLM decontaminator. We apply it to real-world datasets and reveal previously unknown test overlap.
We urge the community to adopt stronger decontamination approaches when using public benchmarks. We call for the community to actively develop fresh one-time exams to accurately evaluate LLMs.

\section*{Acknowledgement}
We would like to express our gratitude to Ying Sheng for the early discussion on rephrased samples.
We also extend our thanks to Dacheng Li, Erran Li, Hao Liu, Jacob Steinhardt, Hao Zhang, and Siyuan Zhuang for providing insightful feedback.
This project is partly supported by gifts from Anyscale, Astronomer, Google, IBM, Intel, Lacework, Microsoft, MBZUAI, Samsung SDS, Uber, and VMware. Lianmin Zheng is supported by a Meta Ph.D. Fellowship.


% 1. any unhidden test set can be contaminated using 
%  rephrase attack

% 2. there is no feasible detection (low cost, accurate)

% 3. unintentional contamination is unavoidable

% 4. one-time exam is a good solution

% \DL{your goal is not to investigate this.}

% \weilin{in this work, we study ...}
% \weilin{you didn't get my point. I meant we need a sentence like: in this work, we study benchmark contamination in the context of large language model... some high-level picture of our work, and then go specific}

% To investigate the capability of current decontamination methods, we study ``rephrased samples'' and propose the corresponding rephrasing algorithm. 
% On the rephrased datasets we create for the MMLU, HumanEval, and GSM-8k tests, we achieve abnormally high scores while successfully evading existing decontamination methods. We verify rephrased samples in open datasets.
% For the MMLU, HumanEval, and GSM-8k tests, the rephrased datasets we craft not only yield dramatically (fake) high scores but also adeptly elude existing decontamination techniques.
% In addition, we present LLM decontaminator that can relaibly detect possible contaminations, e.g. rephased samples.
%To detect possible contamination with rephrased samples, we present the LLM decontaminator.
% Leveraging this method, we identify numerous rephrased samples in public datasets.
% In conclusion, we hope the community will reconsider potential contaminations and utilize the LLM decontaminator to check for rephrased samples while developing datasets. To more precisely evaluate LLM capabilities, the community could take into account strategies like Leetcode's weekly competitions or Kaggle competitions for future benchmarks.
%In conclusion, we urge the community to remain vigilant about possible contamination. 
% We advocate for the adoption of the LLM decontaminator in dataset development and recommend that future benchmarks consider dynamic evaluation mechanisms such as one-time exams.

% Additionally, we introduced a new detection technique called LLM-judge, which outperforms other methods in terms of accuracy. Furthermore, we propose a method to enhance benchmarks called perturbation, which can significantly reduce the likelihood of model cheating. We hope to establish a mechanism similar to Leetcode's weekly competitions, which can accurately reflect the capabilities of models. 

% We identified and defined the phenomenon of rephrase contamination. To explore its impact on benchmark scores, we introduced the rephrase attack to simulate real-world dataset contamination. 
% On the rephrased datasets we created for the MMLU, HumanEval, and GSM-8k tests, we achieved abnormally high scores. We suggested the LLM Decontaminator as a tool for detecting rephrase contamination. This detector is capable of quickly and precisely detecting dataset rephrase contamination. Using the LLM Decontaminator, we found many rephrased test cases in open datasets.
% In conclusion, we hope the community will reconsider potential contaminations and utilize the LLM Decontaminator to check for contamination while developing datasets. To more precisely evaluate LLM capabilities, the community could take into account strategies like Leetcode's weekly challenges or Kaggle competitions for future benchmarks.




